\subsection{伪像测度}

定义 1. --- 设 T 与 B 是两个局部紧空间， \(\mu\) 是 T 上的正测度， p 是 T 到 B 中的 \(\mu\) -可测映射。 B 上的正测度 \(\nu\) 称为 \(\mu\) 在 p 下的伪像测度，如果它满足下列条件： B 的子集 N 局部 \(\nu\) -可忽略，当且仅当 \(\overline{p}^{-1}(N)\) 局部 \(\mu\) -可忽略。

例. --- 1) 若 \(p\) 是 \(\mu\) -逆紧的且 \(\nu = p(\mu)\) ，则 \(\nu\) 是 \(\mu\) 在 \(p\) 下的伪像测度（Ch. V, §6, No.~2, 命题 2 的推论 2）。

\begin{enumerate}
\def\labelenumi{\arabic{enumi})}
\setcounter{enumi}{1}
\tightlist
\item
  设 \(B'\) 是一个局部紧空间， \(\nu'\) 是 \(B'\) 上的正测度；取 T 为空间 \(B \times B'\) ，取 \(\mu\) 为测度 \(\nu \otimes \nu'\) ；若 p 是 T 到 B 上的投影，则 \(\nu\) 是 \(\mu\) 在 p 下的伪像（Ch. V, §8, No.~2, 命题 4 与 No.~3, 命题 7 的推论 1）。
\end{enumerate}

注意，若 \(\nu\) 是 \(\mu\) 在 \(p\) 下的伪像测度，则 \(\nu\) 由 \(p(T)\) 承载。

若 \(\nu\) 是 \(\mu\) 在 p 下的伪像，则在 p 下 \(\mu\) 的伪像测度所成的集合就是与 \(\nu\) 等价的正测度类，且每个与 \(\mu\) 等价的正测度在 p 下有相同的伪像测度。 \(\nu\) 的类称为 \(\mu\) 的类在 p 下的伪像类。

命题 1. --- 设 T 是一个无穷远处可数的局部紧空间， \(\mu\) 是 T 上的正测度， p 是 T 到局部紧空间 B 中的 \(\mu\) -可测映射。则存在 \(\mu\) 在 p 下的一个伪像测度。

因为 T 上存在一个有界测度 \(\mu'\) ，它等价于 \(\mu\) （Ch. V, §5, No.~6, 命题 11）；这时 \(p\) 是 \(\mu'\) -逆紧的。

